\begingroup
\def\xmin{-300} \def\xmax{100} \def\ymin{0} \def\ymax{18}
\def\maxxa{52.5} \def\maxya{11.3} \def\maxxb{93} \def\maxyb{16.5} % CHANGE: read off your graph
\def\minxa{52.5} \def\minya{11.9} \def\minxb{93} \def\minyb{16.1} % CHANGE: read off your graph
\pgfmathsetmacro{\gmax}{(\maxyb-\maxya)/(\maxxb-\maxxa)}
\pgfmathsetmacro{\gmin}{(\minyb-\minya)/(\minxb-\minxa)}
\pgfmathsetmacro{\gmean}{(\gmax+\gmin)/2}
\pgfmathsetmacro{\gunc}{abs(\gmax-\gmin)/2}
\pgfmathsetmacro{\zmax}{\maxxa-\maxya/\gmax}
\pgfmathsetmacro{\zmin}{\minxa-\minya/\gmin}
\pgfmathsetmacro{\zmean}{(\zmax+\zmin)/2}
\pgfmathsetmacro{\zunc}{abs(\zmax-\zmin)/2}

\section{Practical 2: Charles' law}

\subsection{Aim}
To Investigate Charles' Law
\subsection{Variables}
\begin{description}
  \item[Independent:] Temperature
  \item[Dependent:] Length(Volume)
  \item[Control:] Mass, Pressure
\end{description}

\subsection{Introduction}
Measuring the cross-sectional area of the capillary tube is difficult. Since the area stays constant throughout, the length of the fluid in the capillary tube is proportional to the volume of the fluid in the capillary tube:
\begin{equation}
  L \propto V
\end{equation}

\subsection{Method}
\begin{enumerate}
  \item Set up the apparatus as shown in the diagram, ensuring the hot water in the beaker or measuring cylinder fully submerges the air sample.
  \item Stir the water well using the thermometer, then read and record the initial temperature ($\theta_{\text{i}}$ / \unit{\celsius}) and the length of the air sample ($L$ / \unit{\centi\metre}).
  \item Add ice and stir well, taking measurements every \qty{5}{\celsius} down to room temperature.
  \item Plot a graph of $\theta$ / \unit{\celsius} (x-axis) against $L$ / \unit{\centi\metre} (y-axis), drawing a best-fit line.
  \item Deduce the point of absolute zero, $\theta_0$, from the graph by extrapolating the line to the x-intercept (where $L = 0$).
\end{enumerate}

\subsection{Risk assessment}
\begin{table}[H]\centering
\begin{tabularx}{\textwidth}{@{}l X X@{}}
\toprule
Risk & Hazard & Mitigation(s) \\
\midrule
Boiling water & Burns & Handle boiling water with care, taking caution not to cause any spills. \\
\bottomrule
\end{tabularx}
\end{table}

\subsection{Results}

\begin{table}[H]\centering
\pgfplotstabletypeset[
  col sep=comma,
  columns={T,L},
  columns/T/.style={column name={$\theta$ / \unit{\celsius}}, fixed, fixed zerofill, precision=1},
  columns/L/.style={column name={$L$ / \unit{\centi\metre}}, fixed, fixed zerofill, precision=1},
  every head row/.style={before row=\toprule, after row=\midrule},
  every last row/.style={after row=\bottomrule}
]{data-charles.csv}
\caption{Results table.}
\end{table}

\begin{figure}[H]\centering
\begin{tikzpicture}
\begin{axis}[
  width=0.85\textwidth, height=8cm,
  xlabel={$\theta$ / \unit{\celsius}}, ylabel={$L$ / \unit{\centi\metre}},
  xmin=\xmin, xmax=\xmax, ymin=\ymin, ymax=\ymax,
  grid=both, minor tick num=1,
  legend pos=north west,
]

\addplot[only marks, mark=x] table[col sep=comma, x=T, y=L]{data-charles.csv};
\addlegendentry{Data}

\addplot[blue, no marks, forget plot] table[col sep=comma, x=T,
  y={create col/linear regression={y=L}}]{data-charles.csv};
\addplot[blue, no marks, domain=\xmin:\xmax]
  {\pgfplotstableregressiona*x + \pgfplotstableregressionb};
\addlegendentry{Best fit}
\addplot[red, dashed, no marks, domain=\xmin:\xmax] {\gmax*(x-\maxxa)+\maxya};
\addlegendentry{Max gradient}
\addplot[green!50!black, dashed, no marks, domain=\xmin:\xmax] {\gmin*(x-\minxa)+\minya};
\addlegendentry{Min gradient}
\end{axis}
\end{tikzpicture}
\caption{Graph of $L$ against $\theta$}
\end{figure}

\subsection{Analysis}
Gradient of maximum line: $\num[round-mode=places,round-precision=3]{\gmax}$\\
Gradient of minimum line: $\num[round-mode=places,round-precision=3]{\gmin}$\\
Mean gradient: $\num[round-mode=places,round-precision=3]{\gmean}$\\
Uncertainty in gradient: $\pm\,\num[round-mode=places,round-precision=3]{\gunc}$

$\theta_0 = \num[round-mode=places,round-precision=0]{\zmean} \pm
\num[round-mode=places,round-precision=0]{\zunc}\,\unit{\celsius}$

\begin{equation}
  V \propto T, \qquad T\,/\,\unit{\kelvin} = \theta\,/\,\unit{\celsius} + 273.15
\end{equation}

\subsection{Uncertainties}
\subsection{Evaluation}
\endgroup